\documentclass[12pt]{article}

% === amsrefs bibliography setup ===
\usepackage{amsrefs}
\BibSpec{article}{%
  +{}{\PrintAuthors}                {author}
  +{,}{ \textit}                    {title}
  +{.}{ }                           {part}
  +{:}{ }                           {subtitle}
  +{,}{ \PrintContributions}        {contribution}
  +{.}{ \PrintPartials}             {partial}
  +{,}{ }                           {journal}
  +{}{ \textbf}                     {volume}
  +{}{ \parenthesize}               {date}
  +{,}{ }                           {pages}
  +{,}{ }                           {status}
  +{,}{ \PrintDOI}                  {doi}
  +{.}{ Available at \eprint}       {eprint}
  +{}{ \parenthesize}               {transition}
  +{;}{ \PrintReviews}              {review}
  +{.}{}                            {note}
}

% === Core packages ===
\usepackage{amsmath, amsfonts, amssymb, amsthm}
\usepackage{enumerate}
\usepackage{mathrsfs}
\usepackage{xcolor}
\usepackage[colorlinks=true]{hyperref}
\usepackage{framed}
\usepackage{url}
\setlength{\emergencystretch}{3em}

% === Page geometry ===
\setlength{\textwidth}{6.2in}
\setlength{\oddsidemargin}{0.1in}
\setlength{\evensidemargin}{0.1in}
\setlength{\textheight}{8.8in}
\setlength{\topmargin}{0in}

% === Equation numbering ===
\numberwithin{equation}{section}

% === Theorem environments (shared counter) ===
\newtheorem{theorem}{Theorem}[section]
\newtheorem{lemma}[theorem]{Lemma}
\newtheorem{proposition}[theorem]{Proposition}
\newtheorem{corollary}[theorem]{Corollary}
\newtheorem{conjecture}[theorem]{Conjecture}
\theoremstyle{definition}
\newtheorem{definition}[theorem]{Definition}
\newtheorem{hypothesis}[theorem]{Hypothesis}
\newtheorem{remark}[theorem]{Remark}
\newtheorem{goal}[theorem]{Goal}

% === Standard macros ===
\newcommand{\R}{\mathbb{R}}
\newcommand{\N}{\mathbb{N}}
\newcommand{\E}{\mathbb{E}}
\newcommand{\W}{\dot W}
\newcommand{\ud}{\,\mathrm{d}}
\newcommand{\cE}{\mathcal{E}}
\newcommand{\myCmts}[1]{\begin{framed}\noindent\textcolor{red}{#1}\end{framed}}

\begin{document}

\title{The second moment of the hyperbolic Anderson model with time-dependent
  noise: Stratonovich regime\\[0.3em]
  {\normalsize Working notes (Gate G0, 2026-10-27)}}
\author{Le Chen}
\date{\today}
\maketitle

\begin{abstract}
  For the Stratonovich hyperbolic Anderson model driven by a Gaussian noise
  that is fractional in time, X.~Chen~\cite{chen:author2026:time-dependency}
  proves the sharp existence condition and two-sided bounds on the large-time
  growth of the first moment, and he conjectures the growth of all integer
  moments. The source records the second-moment upper bound
  $\E u^2(t,x)\le e^{Ct^2}$. The goal of this project is (i) an upper bound
  for $\E u^2(t,x)$ at the correct scale
  $t^{(4-\alpha-\alpha_0)/(3-\alpha)}$, and (ii) the sharp constant in the
  second-moment asymptotics, or two-sided bounds of the same type as the
  first-moment result. These notes fix the model, record exactly what is
  known (with source pinpoints), state the targets, and list the planned
  route. Section~\ref{S:G0} gives the mixed-term expansion and the upper
  bound $\E u^2(t,x)\le\exp(Ct^\rho)$ for large $t$, where
  $\rho=(4-\alpha-\alpha_0)/(3-\alpha-\alpha_0)$ satisfies
  $\kappa<\rho<2$. It identifies a time-uniformity issue in the source's
  intermediate estimate and recovers the quadratic-exponential bound
  with a corrected argument. The uniform pairing estimate cannot reach
  the target scale. The two principal goals remain open.
\end{abstract}

\tableofcontents

\section{Model and standing assumptions}\label{S:model}

We study the hyperbolic Anderson model
\cite{chen:author2026:time-dependency}*{(1.1)}
\begin{equation}\label{E:HAM}
  \frac{\partial^2 u}{\partial t^2}(t,x)=\Delta u(t,x)+\W(t,x)u(t,x),
  \qquad u(0,x)=1,\quad \frac{\partial u}{\partial t}(0,x)=0,
\end{equation}
for $(t,x)\in\R_+\times\R^d$, $d\in\{1,2,3\}$, interpreted in the
Stratonovich sense (limit of classical wave equations driven by smoothed
noises $\W_{\epsilon,\delta}$). The noise is centered Gaussian with
\begin{equation}\label{E:cov}
  \operatorname{Cov}\big(\W(s,x),\W(t,y)\big)=|s-t|^{-\alpha_0}\gamma(x-y),
  \qquad 0<\alpha_0<1,
\end{equation}
where $\gamma\ge0$ is nonnegative definite with spectral measure $\mu$.

\begin{hypothesis}[Homogeneity]\label{H:homog}
  $\gamma(cx)=c^{-\alpha}\gamma(x)$ for $c>0$, and
  \begin{equation}\label{E:threshold}
    \alpha_0+\alpha<2,
  \end{equation}
  with $0<\alpha<d$, or $\alpha=d=1$ (spatial white noise).
\end{hypothesis}

Throughout we write
\begin{equation}\label{E:kappa}
  \kappa:=\frac{4-\alpha-\alpha_0}{3-\alpha}\in(1,2).
\end{equation}

\section{What is known}\label{S:known}

All pinpoints below refer to the current author version
\cite{chen:author2026:time-dependency} (64 pages, ``Submitted to the Annals
of Applied Probability'', SHA-256 \texttt{5f7385ca\dots}); see
\path{sources/SOURCE-MANIFEST.md}. arXiv:2510.01412v1
\cite{chen:25:time-dependency} differs in Theorem~1.2 and must not be used as
the authority for it.

\subsection{Stratonovich expansion and moments}

The solution is the Stratonovich expansion
\cite{chen:author2026:time-dependency}*{(2.13)--(2.16)}
\begin{equation}\label{E:expansion}
  \begin{split}
    u(t,x)&=\sum_{n=0}^\infty S_n\big(g_n(\cdot,t,x)\big),\\
    g_n(s_1,\dots,s_n,x_1,\dots,x_n,t,x)
    &=\Big(\prod_{l=1}^nG(s_{l+1}-s_l,x_{l+1}-x_l)\Big)
    \mathbf 1_{[0,t]^n_<}(s_1,\dots,s_n),
  \end{split}
\end{equation}
with $s_{n+1}=t$, $x_{n+1}=x$, $G$ the fundamental solution of the wave
equation, and $[0,t]^n_<=\{0<s_1<\dots<s_n<t\}$. Odd expected Stratonovich
integrals vanish and even ones are given by Wick's formula
\cite{chen:author2026:time-dependency}*{(2.19)--(2.21)}, a sum over pair
partitions $\mathcal D\in\Pi_n$ of $\{1,\dots,2n\}$ that includes
self-pairings within one chain. The second moment of a single chaos term is
the two-chain pairing integral \cite{chen:author2026:time-dependency}*{Theorem~4.1, (4.3)}:
\begin{equation}\label{E:Sn2}
  \begin{split}
    \E\big[S_n(g_n(\cdot,t,x))\big]^2
    &=\sum_{\mathcal D\in\Pi_n}\int_{[0,t]^n_<\times[0,t]^n_<}\ud s
    \int_{(\R^d)^n\times(\R^d)^n}\ud x
    \Big(\prod_{\rho=1}^2\prod_{l\in I_\rho}G(s_l-s_{l-1},x_l-x_{l-1})\Big)\\
    &\qquad\times\prod_{(j,k)\in\mathcal D}|s_j-s_k|^{-\alpha_0}\gamma(x_j-x_k),
  \end{split}
\end{equation}
with $I_1=\{1,\dots,n\}$, $I_2=\{n+1,\dots,2n\}$ and the conventions of
\cite{chen:author2026:time-dependency}*{Theorem~4.1}.

The full mixed-term expansion is justified in Proposition~\ref{P:mixed}.

\subsection{The universal two-chain bound}

\cite{chen:author2026:time-dependency}*{Theorem~3.1, (3.6), p.~15} bounds
exactly such two-chain pairing integrals with general nonnegative kernels
$G_l$ by products of the norms $\|G_l\|^{(0)}$, $\|G_l\|^{(1)}_{j,k}$,
$\|G_l\|^{(2)}_{j,k}$ of \cite{chen:author2026:time-dependency}*{(3.3)--(3.5)},
over a partition $\{Q_0,Q_1,Q_2\}$ of $\{1,\dots,2n\}$ with $\#Q_0=\#Q_1$ and
$\#Q_2$ even. As stated, the theorem uses finite spatial spectral measures
$\mu_{j,k}$ with nonnegative covariance functions and temporal kernels
$\Gamma(\alpha_0)^{-1}\int_{A_{j,k}}e^{-\lambda|r|}
\lambda^{\alpha_0-1}\ud\lambda$ for measurable $A_{j,k}\subset\R_+$.
The application to the full covariance therefore requires regularization;
this is carried out in Lemma~\ref{L:crude-corrected}.
The theorem includes cross-pairings between the two chains. In the
source it is used only for Stratonovich integrability (Section~4) and, via
Lemma~A.1 with $G_l(t,x)=e^{-\theta t}G(t,x)$, in Laplace variables.

\subsection{Known moment bounds}

\begin{theorem}[\cite{chen:author2026:time-dependency}*{Theorem~1.2, (1.10)--(1.13)}]\label{T:first}
  Under Hypothesis~\ref{H:homog},
  \[
    L_*\le\liminf_{t\to\infty}t^{-\kappa}\log\E u(t,x)
    \le\limsup_{t\to\infty}t^{-\kappa}\log\E u(t,x)\le L^*,
  \]
  where, for $\cE^\sharp\in\{\cE_*,\cE\}$,
  \[
    L(\cE^\sharp):=(3-\alpha)\bigg(
    \frac{(4-\alpha-2\alpha_0)^{\frac{4-\alpha-2\alpha_0}{2}}}
    {2^3(4-\alpha-\alpha_0)^{4-\alpha-\alpha_0}}
    \Big(\frac{2\cE^\sharp}{2-\alpha}\Big)^{\frac{2-\alpha}{2}}
    \bigg)^{\frac1{3-\alpha}},
    \qquad L_*=L(\cE_*),\quad L^*=L(\cE),
  \]
  and $\cE_*<\cE$ are the variational quantities
  \cite{chen:author2026:time-dependency}*{(1.11)--(1.13)}.
\end{theorem}

The source predicts that the limit exists and equals $L^*$
\cite{chen:author2026:time-dependency}*{p.~4}; this is not proved, and
\cite{chen:author2026:time-dependency}*{Remark~8.1} explains why the lower
bound is not sharp (complex-valued intersection local time).

\begin{theorem}[\cite{chen:author2026:time-dependency}*{Remark~5.1, (5.8)}]\label{T:crude}
  $\E u^2(t,x)\le e^{Ct^2}$ for large $t$.
\end{theorem}

\begin{conjecture}[\cite{chen:author2026:time-dependency}*{(1.22)}]\label{Conj:allp}
  For $p=1,2,\dots$,
  \[
    \lim_{t\to\infty}t^{-\kappa}\log\E|u(t,x)|^p
    =p^{\frac{4-\alpha}{3-\alpha}}L^*.
  \]
\end{conjecture}

Since $p^{(4-\alpha)/(3-\alpha)}$ grows faster than $p$, the conjectured
second-moment constant $2^{(4-\alpha)/(3-\alpha)}L^*$ exceeds $2L^*$.
The source explains why its method stops at $p=1$
\cite{chen:author2026:time-dependency}*{Remark~6.3}: the representation of
Theorem~6.1 there uses the monotone order $s_1\le\dots\le s_{2n}$ of a single
chain.

\begin{proposition}[Jensen lower bound]\label{P:Jensen}
  Under Hypothesis~\ref{H:homog},
  \[
    \liminf_{t\to\infty}t^{-\kappa}\log\E u^2(t,x)\ge 2L_*.
  \]
\end{proposition}

\begin{proof}
  By \eqref{E:expansion} and \cite{chen:author2026:time-dependency}*{(2.20)--(2.21)},
  $\E u(t,x)=\sum_{n\ge0}\E S_{2n}(g_{2n}(\cdot,t,x))\ge1$, since every term is
  a pairing integral of a nonnegative integrand and $S_0=1$. By the
  Cauchy--Schwarz inequality, $\E u^2(t,x)\ge(\E u(t,x))^2$, so
  $t^{-\kappa}\log\E u^2(t,x)\ge2\,t^{-\kappa}\log\E u(t,x)$, and the claim
  follows from the lower bound in Theorem~\ref{T:first}. This proves Proposition~\ref{P:Jensen}.
\end{proof}

\section{Gate G0: cross terms and the loss in the norm bound}\label{S:G0}

This section records a diagnosis of the proposed route, not a proof of
Goal~\ref{G:upper}. Put $X_n(t,x)=S_n(g_n(\cdot,t,x))$, with $X_0=1$, and
\[
  \beta=4-\alpha-\alpha_0>2,\qquad
  P_k=\#\Pi_k=\frac{(2k)!}{2^k k!}.
\]
The count follows by pairing the smallest remaining label with one of
$2k-1$ labels: $P_0=1$ and $P_k=(2k-1)P_{k-1}$, which gives the displayed
formula by induction. All constants below may depend on the fixed covariance and its parameters,
but not on the orders or on $t$.

\subsection{The complete two-chain expansion}

For $m+n=2k$, label the two chains by
$I_1=\{1,\ldots,m\}$ and $I_2=\{m+1,\ldots,m+n\}$.
Let $p(l)$ denote the preceding vertex in its own chain; at each first
vertex set $(s_{p(l)},y_{p(l)})=(0,0)$. Thus, the second chain has its own
root, rather than the last vertex of the first chain. For $D\in\Pi_k$, set
\begin{equation}\label{E:mixed-diagram}
  \begin{split}
    J_{m,n,D}(t)
    ={}&\int_{[0,t]^m_<\times[0,t]^n_<}\ud\mathbf s
    \int_{(\R^d)^{m+n}}\ud\mathbf y\,
    \prod_{\rho=1}^2\prod_{l\in I_\rho}
    G(s_l-s_{p(l)},y_l-y_{p(l)})\\
    &\hspace{20mm}\times
    \prod_{(j,h)\in D}|s_j-s_h|^{-\alpha_0}\gamma(y_j-y_h).
  \end{split}
\end{equation}
An empty chain contributes the factor one. For $d=3$, the spatial Green
kernels are measures; for spatial white noise the covariance is a delta
measure. Integrals in~\eqref{E:mixed-diagram} use the generalized-kernel
and regularization conventions of the source, not pointwise values of
these measures; see~\cite{chen:author2026:time-dependency}*{Remark~2.3,
pp.~13--14, and Theorem~4.1}.

\begin{proposition}[Mixed terms]\label{P:mixed}
  Under Hypothesis~\ref{H:homog},
  \begin{equation}\label{E:mixed-moment}
    \E[X_m(t,x)X_n(t,x)]
    =\begin{cases}
      \displaystyle\sum_{D\in\Pi_k}J_{m,n,D}(t),&m+n=2k,\\
      0,&m+n\text{ odd},
    \end{cases}
  \end{equation}
  where the case $m=n=0$ equals one. In particular,
  \begin{equation}\label{E:full-second}
    \E u^2(t,x)=\sum_{k=0}^\infty\sum_{m=0}^{2k}
    \sum_{D\in\Pi_k}J_{m,2k-m,D}(t).
  \end{equation}
  This is a convergent series of nonnegative terms.
\end{proposition}

\begin{proof}
  From~\cite{chen:author2026:time-dependency}*{Theorem~4.1}, each
  deterministic kernel $g_n$ is Stratonovich integrable. Its spectral
  hypothesis (1.5) follows from Hypothesis~\ref{H:homog}, as recorded
  on p.~3 of the source. Lemma~2.5 and
  (2.18) of that source apply to these deterministic kernels and give
  $X_mX_n=S_{m+n}(g_m\otimes g_n)$ for $m,n\ge1$. More explicitly,
  the smoothed factors converge in $L^4$, so their products converge in
  $L^2$ by H\"older's inequality. Thus, the product rule does not require
  an unproved interchange of infinite series. The cases with a zero order
  use $X_0=1$. From the source's (2.20)--(2.21), the expectation is zero
  at odd total order and is the Wick sum at even total order.

  We also justify the individual generalized integrals in the Wick sum.
  Fix $D\in\Pi_k$ and take $k$ independent copies $W^e$, $e\in D$, of
  the centered Gaussian noise $W$, each with the same covariance as $W$.
  At vertex $l$ use the copy indexed by the unique pair $e\in D$
  containing $l$. For either nonempty chain of length $r$, let
  $Y_{r,a}$ be its integral against the product of these smoothed noises,
  where $a$ denotes all smoothing parameters on that chain. For two
  parameter choices $a,b$, Wick's formula expresses
  $\E[Y_{r,a}Y_{r,b}]$ as a finite sum over those pairings of two copies
  of this length-$r$ chain that join vertices of the same color.
  This is a subset of the equal-length diagrams in
  \cite{chen:author2026:time-dependency}*{Lemma~4.4, p.~30}.
  Each summand has a finite joint limit as $a,b\to0$, independent of
  their relative rates. Consequently, the three terms in
  \[
    \E|Y_{r,a}-Y_{r,b}|^2
    =\E Y_{r,a}^2+\E Y_{r,b}^2-2\E[Y_{r,a}Y_{r,b}]
  \]
  have the same limiting covariance, and $Y_{r,a}$ converges in $L^2$.
  For $d=3$, these integrals use precisely the measure-valued Green-kernel
  convention of the source's pp.~13--14 and Lemma~4.4: for fixed noise
  smoothing first take the $G_\eta$ limit as $\eta\downarrow0$, and then
  take the noise-smoothing limit just proved.
  An empty chain is the constant one. The product of the two colored
  chain integrals therefore converges in $L^1$, by Cauchy--Schwarz.
  In that product each color occurs exactly twice in total, so
  independence and Wick's formula leave exactly the prescribed pairing
  $D$. Its expectation is the regularized integral defining
  $J_{m,n,D}(t)$, after the same reversal of both chains used below.
  Thus, each such integral has a finite limit. The regularized
  integrands are nonnegative, so $J_{m,n,D}(t)\ge0$.

  In each chain reverse the order of its variables, replace each time by
  $t$ minus that time, and each spatial coordinate by $x$ minus that
  coordinate. Each terminal Green factor becomes an initial Green factor.
  Both time and spatial covariance differences change sign; the covariance
  kernels are even. The resulting permutation of the $2k$ labels is a
  bijection of $\Pi_k$. This gives~\eqref{E:mixed-moment}, including its
  nonnegativity. These are exactly the two chains in the integrand of
  the source's (3.6), with $n_1=m$, $n_2=n$, restricted to the finite
  horizons. The same transformations apply to the source's regularized
  integrals before taking their defining limits.

  From~\cite{chen:author2026:time-dependency}*{Theorem~1.1 and Section~5,
  p.~34}, $u$ is the solution represented by the series of the $X_n$.
  Its fixed-$t$ summability is also proved independently in
  Lemma~\ref{L:crude-corrected} below, whose proof uses only the source's
  diagonal identity and Theorem~3.1, not the present mixed expansion.
  In particular, $U_N=\sum_{n=0}^N X_n$ converges to $u$ in $L^2$.
  Consequently,
  \[
    |\E U_N^2-\E u^2|
    \le\|U_N-u\|_2(\|U_N\|_2+\|u\|_2)\longrightarrow0.
  \]
  The finite expansion of $\E U_N^2$ is a sum of the nonnegative mixed
  expectations already computed. Its increasing limit is therefore the
  full double sum, independently of its enumeration. Regrouping by total
  order gives~\eqref{E:full-second}. This proves Proposition~\ref{P:mixed}.
\end{proof}

\subsection{A small-order check and a source uniformity issue}

\begin{lemma}[White-space test]\label{L:white-test}
  For $d=\alpha=1$ and $\gamma=\delta_0$, put
  $c_0=((1-\alpha_0)(2-\alpha_0)(3-\alpha_0))^{-1}$. Then
  \begin{equation}\label{E:white-test}
    \E X_1(t,x)^2=c_0t^{3-\alpha_0},\qquad
    \E X_2(t,x)=\tfrac12c_0t^{3-\alpha_0}.
  \end{equation}
  In particular, the total-order-two contribution to~\eqref{E:full-second}
  is $2c_0t^{3-\alpha_0}$.
\end{lemma}

\begin{proof}
  Here $G(s,y)=\frac12\mathbf1_{\{|y|<s\}}$. Therefore,
  $\int_\R G(s,y)G(r,y)\ud y=\min(s,r)/2$. The $(1,1)$ integral is
  \[
    \int_0^t\ud s\int_0^s r(s-r)^{-\alpha_0}\ud r
    =\frac{t^{3-\alpha_0}}
    {(1-\alpha_0)(2-\alpha_0)(3-\alpha_0)}.
  \]
  For the $(0,2)$ integral, the delta covariance sets the two spatial
  vertices equal. The integral of the first Green factor is $r$, whereas
  the second Green factor at zero is $1/2$. Its time integral is half the
  preceding display. The $(2,0)$ term is identical. This proves
  Lemma~\ref{L:white-test}.
\end{proof}

\begin{remark}[Source fidelity]\label{R:source-uniformity}
  The author version's Remark~5.1 refers to (4.4), but (4.4), on p.~22,
  is a smoothing limit. The relevant bound is (5.4), on p.~35,
  $\E X_n(t,x)^2\le C^nt^{2n}/n!$, with the preceding sentence asserting
  that $C$ is independent of $n$ and $t$. Taken for all $t>0$, this
  uniformity contradicts~\eqref{E:white-test} at $n=1$.
  The proof on pp.~34--35 sets $\theta=n/t$. Its first-norm estimate
  uses (A.1), which is finite for every positive $\theta$ but does not
  give a constant uniform as $\theta\downarrow0$. Its second-norm
  estimate uses the large-$\theta$ bound (A.3) without retaining that
  restriction. Both losses of uniformity matter when $n/t$ is small.
  This is a source-internal uniformity discrepancy; it does not refute
  Remark~5.1's final large-time conclusion. The next lemma repairs this
  particular deduction without importing the asserted uniformity of (5.4).
\end{remark}

\subsection{Recovering the quadratic-exponential bound}

Write $G_\theta(s,y)=e^{-\theta s}G(s,y)$ and let
$N_i(\theta)$, $i=0,1,2$, denote the three norms in the source's
(3.3)--(3.5), with the full covariance. From its (2.4) and Lemma~A.1,
there are $C_0<\infty$ and $\theta_0>0$ such that, for
$\theta\ge\theta_0$,
\begin{equation}\label{E:crude-norms}
  N_0(\theta)=\theta^{-2},\qquad
  N_1(\theta)\le C_0,\qquad N_2(\theta)\le C_0\theta^{-1}.
\end{equation}
For $N_1$, the spectral formula (A.1) is decreasing in $\theta$, so its
finite value at $\theta_0$ suffices. The last bound is (A.3).

\begin{lemma}[Corrected crude estimate]\label{L:crude-corrected}
  Under Hypothesis~\ref{H:homog}, for suitable $C<\infty$ and
  $\theta_0>0$,
  \begin{equation}\label{E:crude-corrected}
    \E X_n(t,x)^2\le e^{2\theta_0t}\frac{C^nt^{2n}}{n!}
    \quad(n\ge1,t>0),\qquad
    \E u^2(t,x)\le2e^{2\theta_0t+2Ct^2}.
  \end{equation}
\end{lemma}

\begin{proof}
  On the two finite chains, the product of the exponential weights
  telescopes to $e^{-\theta(s_n+s_{2n})}\ge e^{-2\theta t}$.
  Extend the nonnegative integral to the two infinite time simplices
  and apply the source's Theorem~3.1. If $q_i=\#Q_i$, then
  $q_0=q_1$ and $q_0+q_2/2=n$. Hence,~\eqref{E:crude-norms} bounds
  each pairing by $e^{2\theta t}C^n\theta^{-2n}$.
  Summing the $P_n$ pairings gives
  \begin{equation}\label{E:theta-envelope}
    \E X_n(t,x)^2\le P_n e^{2\theta t}C^n\theta^{-2n},
    \qquad\theta\ge\theta_0.
  \end{equation}
  To be precise about the source's finite-spectral-measure hypothesis in
  Theorem~3.1, first apply it with the spatially smoothed covariance and
  the temporal truncations of (2.24). The smoothed spectral measure is
  finite: the Gaussian multiplier is bounded by a constant times
  $(1+|\xi|^2)^{-(2-\alpha_0)/2}$, whose integral is finite by (1.5).
  This argument requires no radial density for the spectral measure.
  The zero norm is unchanged; the
  first norm is bounded by (A.1), since its spectral multiplier is at
  most one; the squared second norm is bounded by the full temporal
  covariance norm and then by removal of the spatial multiplier in the
  nonnegative space-time spectral integral, as in (4.14). The bound is
  uniform in these regularizations. In $d=3$, first use
  $G_\eta=p_\eta*G$, as defined on pp.~13--14 of the source. Its
  spatial Fourier multiplier is bounded by one and its spatial mass
  equals that of $G$, so the same bounds are uniform in $\eta$ as well.
  For each fixed noise regularization $(\epsilon,\delta)$, first let
  $\eta\downarrow0$ in the source's $L^2$ definition. The squared
  moments converge, so the inequality survives this first limit.
  Then let $(\epsilon,\delta)\to0$ using Theorem~4.1. In $d=1,2$ only
  this latter limit is needed. These limits transfer a uniform inequality;
  no pointwise delta covariance bound is being used.

  Choose $\theta=\theta_0+n/t$ in~\eqref{E:theta-envelope}. Since
  $(2n)!\le(2n)^{2n}$, we have $P_n n^{-2n}\le2^n/n!$.
  Absorb $2^ne^{2n}$ in $C^n$ to obtain the first bound
  in~\eqref{E:crude-corrected}. Its square roots are summable for each
  fixed $t$, by the ratio test. Thus, the norm series converges independently
  of Proposition~\ref{P:mixed}. By Minkowski and Cauchy--Schwarz,
  \[
    \begin{split}
      \|u(t,x)\|_2
      &\le e^{\theta_0t}\sum_{n\ge0}\frac{(\sqrt C t)^n}{\sqrt{n!}},\\
      \left(\sum_{n\ge0}\frac{z^n}{\sqrt{n!}}\right)^2
      &\le\left(\sum_{n\ge0}2^{-n}\right)
      \left(\sum_{n\ge0}\frac{(2z^2)^n}{n!}\right)
      =2e^{2z^2}.
    \end{split}
  \]
  This proves Lemma~\ref{L:crude-corrected}.
\end{proof}

Thus, the exponent two enters through replacing the paired norm products
by $C\theta^{-2}$, followed by the pairing count and the order summation.
The sum over unequal chain lengths is not responsible for this loss:
the same loss already occurs in the diagonal estimate and Minkowski's
inequality.

\subsection{What homogeneity recovers, and what the bound still loses}

\begin{proposition}[The uniform pairing estimate]\label{P:envelope}
  Under Hypothesis~\ref{H:homog}, there is $K>0$ such that
  \begin{equation}\label{E:homogeneous-envelope}
    \E[X_m(t,x)X_n(t,x)]
    \le P_k K^k e^{2\theta t}\theta^{-\beta k},
    \qquad m+n=2k\ge2,\quad\theta>0.
  \end{equation}
  Define the numerical majorant produced by applying this estimate
  separately at every total order by
  \begin{equation}\label{E:B-envelope}
    B_K(t)=1+\sum_{k\ge1}(2k+1)P_kK^k
    \inf_{\theta>0}\{e^{2\theta t}\theta^{-\beta k}\}.
  \end{equation}
  Then $\E u^2(t,x)\le B_K(t)$ and, for some $c,C>0$ and all large $t$,
  \begin{equation}\label{E:envelope-growth}
    ct^{\rho}\le\log B_K(t)\le Ct^{\rho},\qquad
    \rho=\frac{\beta}{\beta-1}
    =\frac{4-\alpha-\alpha_0}{3-\alpha-\alpha_0}>\kappa.
  \end{equation}
  The lower bound here concerns the numerical majorant $B_K$, not the
  second moment.
\end{proposition}

\begin{proof}
  The wave kernel scales as
  $G(cs,c\ud y)=cG(s,\ud y)$, and the covariance has degree $-\alpha$.
  Changing all time and spatial variables by $\theta^{-1}$ in the norms
  yields
  \begin{equation}\label{E:homogeneous-norms}
    N_0(\theta)=\theta^{-2},\quad
    N_1(\theta)=N_1(1)\theta^{\alpha+\alpha_0-2},\quad
    N_2(\theta)^2=N_2(1)^2\theta^{-\beta}.
  \end{equation}
  For the first norm this also follows directly from (A.1): substitute
  $\lambda=\theta v$, $\xi=\theta\eta$, using the degree $\alpha$ of
  the spectral measure. All constants are finite by Lemma~A.1.
  The products $2N_0N_1$ and $N_2^2$ therefore both have degree
  $-\beta$. Choose
  $K\ge\max\{2N_1(1),N_2(1)^2\}$, increasing it to a positive
  constant if needed. Theorem~3.1 and $q_0+q_2/2=k$ give
  $K^k\theta^{-\beta k}$ for every partition it supplies. The finite
  horizons contribute at most $e^{2\theta t}$. Apply this first to
  regularized moments. In $d=3$, fix $(\epsilon,\delta)$ and let
  $\eta\downarrow0$ in each factor's $L^2$ limit; Cauchy--Schwarz gives
  $L^1$ convergence of their product. Then let
  $(\epsilon,\delta)\to0$ in the same way, using Theorem~4.1 for each
  chain. In $d=1,2$ only the noise limit is required. The estimates are
  uniform at both stages, by the multiplier bounds in
  Lemma~\ref{L:crude-corrected}. This proves~\eqref{E:homogeneous-envelope}.
  Empty-chain cases satisfy the same bound (using $e^{2\theta t}$ is
  harmless). Proposition~\ref{P:mixed} now gives the majorant.

  The infimum in~\eqref{E:B-envelope} is attained at
  $\theta=\beta k/(2t)$ and equals
  $e^{\beta k}(2t/(\beta k))^{\beta k}$. Put $q=\beta-1>1$.
  The elementary estimates
  \[
    k!\le P_k\le2^k k!,\qquad
    (k/e)^k\le k!\le k^k,\qquad 2k+1\le3^k
  \]
  show that, for fixed positive $a,b$ depending only on $K,\beta$,
  \[
    \sum_{k\ge0}\frac{(at^\beta)^k}{(k!)^q}
    \le B_K(t)\le
    \sum_{k\ge0}\frac{(bt^\beta)^k}{(k!)^q}.
  \]
  Indeed, the additional factor $k!^\beta/k^{\beta k}$ lies between
  $e^{-\beta k}$ and one and can be absorbed into $a^k,b^k$.
  For $z>0$, the right-hand series is at most
  $(\sum_{k\ge0}z^{k/q}/k!)^q=e^{qz^{1/q}}$, by the elementary
  $\ell^q$ versus $\ell^1$ inequality for nonnegative sequences.
  For a lower bound choose $k=\lfloor z^{1/q}/2\rfloor$ and use
  $k!\le k^k$; its single summand is at least $2^{qk}$.
  These estimates prove~\eqref{E:envelope-growth}. Finally,
  $\beta-1=3-\alpha-\alpha_0<3-\alpha$, so $\rho>\kappa$.
  This proves Proposition~\ref{P:envelope}.
\end{proof}

\begin{remark}[G0 verdict and its scope]\label{R:G0-verdict}
  Since $\beta>2$, we have $1<\rho<2$. Thus,
  Proposition~\ref{P:envelope} improves the crude quadratic time scale
  under homogeneity, while leaving Goal~\ref{G:upper} open.
  Theorem~3.1 used through its uniform pairing estimate is \emph{not
  viable for Goal~\ref{G:upper}}: even with exact homogeneous norm
  scaling, the resulting majorant has exponent $\rho>\kappa$.
  The order coefficient has factorial denominator $(k!)^{3-\alpha-\alpha_0}$,
  whereas the desired scale requires $(k!)^{3-\alpha}$; the missing
  factor is $(k!)^{\alpha_0}$, up to exponential factors in $k$.
  Counting the $2k+1$ chain-length splits does not cause this loss,
  since it can be absorbed into an exponential factor.
  This conclusion does not rule out a more refined application retaining
  temporal order or grouping pairings before estimating them, and does not
  disprove the conjecture. In the formal boundary limit
  $\alpha_0\downarrow0$, both $\rho$ and $\kappa$ tend to
  $(4-\alpha)/(3-\alpha)$. This necessary time-independent calibration
  does not distinguish them and asserts no endpoint theorem or uniformity
  of constants in $\alpha_0$. At $p=1$ use the empty-chain case with
  $e^{\theta t}$ in place of $e^{2\theta t}$; minimizing at
  $\theta=\beta k/t$ changes only a constant to the power $k$, leaving
  the same factorial denominator $(k!)^{3-\alpha-\alpha_0}$.
  It therefore also loses the factorial power $\alpha_0$ relative to
  Chen's Theorem~1.2, so the loss is not specific to the second moment. Both the crude and homogeneous upper bounds are
  consistent with the Jensen lower bound, being at larger time scales.
  No identification with the Skorohod second moment is used: its exclusion
  of self-pairings changes the combinatorics, already visible in the
  $(0,2)$ and $(2,0)$ terms of Lemma~\ref{L:white-test}.
\end{remark}

\section{Targets}\label{S:targets}

\begin{goal}[Upper bound at the correct scale]\label{G:upper}
  Under Hypothesis~\ref{H:homog}, prove
  $\limsup_{t\to\infty}t^{-\kappa}\log\E u^2(t,x)<\infty$, ideally with an
  explicit constant.
\end{goal}

\begin{goal}[Sharp second-moment constant]\label{G:sharp}
  Prove two-sided bounds of the type of Theorem~\ref{T:first} for $p=2$:
  \[
    2^{\frac{4-\alpha}{3-\alpha}}L_*
    \le\liminf_{t\to\infty}t^{-\kappa}\log\E u^2(t,x)
    \le\limsup_{t\to\infty}t^{-\kappa}\log\E u^2(t,x)
    \le2^{\frac{4-\alpha}{3-\alpha}}L^*,
  \]
  which would confirm Conjecture~\ref{Conj:allp} at $p=2$ up to the same gap
  $\cE_*<\cE$ as for $p=1$. The exact limit
  $2^{(4-\alpha)/(3-\alpha)}L^*$ is the conditional target, contingent on
  closing that gap.
\end{goal}

\myCmts{The exponents in Goal~\ref{G:sharp} are the conjectured ones. If
  the analysis produces a different $p=2$ constant, that is a finding about
  Conjecture~\ref{Conj:allp}, not a reason to adjust the target.}

\section{Planned route}\label{S:route}

\begin{enumerate}[G0.]
  \item \textbf{Diagnosis (falsification first).} Reproduce
  \item \textbf{Calibrations.} Any proposed bound must (a) as $\alpha_0\to0$
  \item \textbf{Upper bound (Goal~\ref{G:upper}).} Laplace transform in $t$,
  \item \textbf{Two-copy representation (Goal~\ref{G:sharp}).} Each of the
\end{enumerate}
% [held back: collaborators' unpublished work]

\bibliography{All,extra}

\end{document}
